% ==========================================
% 二维连续型随机变量综合应用
% ==========================================

\vspace{1.5em}
\subsubsection*{1. 二维随机变量的边缘与条件密度}

\begin{例题}[求解常数、边缘密度与条件密度]
	已知二维随机变量 $(X,Y)$ 的概率密度函数为
	$$ f(x,y) = \begin{cases} \frac{1}{A}, & -1 \le x \le 1, 0 \le y \le x^2 \\ 0, & \text{其他} \end{cases} $$
	其中 $A$ 为常数。
	
	(i) 试求常数 $A$；
	
	(ii) 求边缘概率密度 $f_Y(y)$；
	
	(iii) 求条件概率密度 $f_{X|Y}(x|y)$。
\end{例题}

\textbf{解答}\quad

\textbf{(i) 求常数 $A$} \\
根据二维概率密度函数的归一性性质，联合密度在整个平面的二重积分必须等于 1，即 $\iint_{-\infty}^{+\infty} f(x,y) dx dy = 1$。\\
积分区域 $D$ 为 $-1 \le x \le 1, 0 \le y \le x^2$。我们选择先对 $y$ 积分，后对 $x$ 积分：
$$ 
\begin{aligned}
\iint_D \frac{1}{A} dx dy &= \frac{1}{A} \int_{-1}^{1} \left( \int_{0}^{x^2} dy \right) dx \\
&= \frac{1}{A} \int_{-1}^{1} x^2 dx \\
&= \frac{1}{A} \left[ \frac{1}{3}x^3 \right]_{-1}^{1} \\
&= \frac{1}{A} \left( \frac{1}{3} - \left(-\frac{1}{3}\right) \right) = \frac{2}{3A}
\end{aligned}
$$
令 $\frac{2}{3A} = 1$，解得 $A = \frac{2}{3}$。\\
因此，联合概率密度函数为：
$$ f(x,y) = \begin{cases} \frac{3}{2}, & -1 \le x \le 1, 0 \le y \le x^2 \\ 0, & \text{其他} \end{cases} $$

\vspace{0.5em}
\textbf{(ii) 求边缘概率密度 $f_Y(y)$} \\
由联合密度的非零区域 $0 \le y \le x^2 \le 1$ 可知，$Y$ 的可能取值范围是 $0 < y < 1$。\\
对于固定的 $y \in (0,1)$，要求原密度函数非零，对应的 $x$ 必须满足 $x^2 \ge y$ 且 $-1 \le x \le 1$。解这个不等式，得到 $x$ 的范围是由两段区间组成的：$x \in [-1, -\sqrt{y}] \cup [\sqrt{y}, 1]$。\\
因此，对 $x$ 积分时需要分段：
$$ 
\begin{aligned}
f_Y(y) &= \int_{-\infty}^{+\infty} f(x,y) dx \\
&= \int_{-1}^{-\sqrt{y}} \frac{3}{2} dx + \int_{\sqrt{y}}^{1} \frac{3}{2} dx \\
&= \frac{3}{2}(-\sqrt{y} - (-1)) + \frac{3}{2}(1 - \sqrt{y}) \\
&= \frac{3}{2}(1 - \sqrt{y}) + \frac{3}{2}(1 - \sqrt{y}) \\
&= 3(1 - \sqrt{y})
\end{aligned}
$$
当 $y$ 不在 $(0, 1)$ 范围内时，$f_Y(y) = 0$。所以边缘概率密度为：
$$ f_Y(y) = \begin{cases} 3(1 - \sqrt{y}), & 0 < y < 1 \\ 0, & \text{其他} \end{cases} $$

\vspace{0.5em}
\textbf{(iii) 求条件概率密度 $f_{X|Y}(x|y)$} \\
根据条件概率密度的定义：$f_{X|Y}(x|y) = \frac{f(x,y)}{f_Y(y)}$（当 $f_Y(y) > 0$ 时）。\\
对于给定的 $0 < y < 1$，有 $f_Y(y) = 3(1 - \sqrt{y})$。将 $f(x,y)$ 和 $f_Y(y)$ 直接代入公式：
$$ f_{X|Y}(x|y) = \frac{3/2}{3(1 - \sqrt{y})} = \frac{1}{2(1 - \sqrt{y})} $$
注意，此条件密度有效的 $x$ 范围即为原联合密度非零的区域（因为给定 $y$ 的前提下，$x$ 只能在特定的范围内取值），即 $x \in [-1, -\sqrt{y}] \cup [\sqrt{y}, 1]$。\\
因此，条件概率密度为：
$$ f_{X|Y}(x|y) = \begin{cases} \frac{1}{2(1 - \sqrt{y})}, & x \in [-1, -\sqrt{y}] \cup [\sqrt{y}, 1] \\ 0, & \text{其他} \end{cases} $$